\documentclass[a4paper,12pt]{article}
\usepackage[a4paper, margin=2.5cm]{geometry}
\usepackage[utf8]{inputenc}
\usepackage[T1]{fontenc}
\usepackage[indonesian]{babel}
\usepackage{booktabs}
\usepackage{longtable}
\usepackage{array}
\usepackage{multirow}
\usepackage{xcolor}
\usepackage{enumitem}

\title{\textbf{SILABUS MATA KULIAH}\\
\textbf{PEMROGRAMAN GAME}}
\author{Program Studi S1 Informatika}
\date{}

\begin{document}

\maketitle

\vspace{1cm}

\begin{table}[h]
\centering
\begin{tabular}{ll}
\textbf{Nama Mata Kuliah} & : Pemrograman Game \\
\textbf{Kode Mata Kuliah} & : IF-XXX \\
\textbf{Bobot SKS} & : 3 SKS \\
\textbf{Semester} & : [Tentukan Semester] \\
\textbf{Program Studi} & : S1 Informatika \\
\textbf{Dosen Pengampu} & : [Nama Dosen] \\
\end{tabular}
\end{table}

\newpage

\section{Deskripsi Mata Kuliah}

Mata kuliah Pemrograman Game merupakan mata kuliah yang membekali mahasiswa dengan pengetahuan dan keterampilan praktis dalam pengembangan game menggunakan Unity Engine dan bahasa pemrograman C\#. Mata kuliah ini mengadopsi pendekatan Outcome-Based Education (OBE) dengan fokus pada pembelajaran berbasis proyek (project-based learning) yang dominan praktik dibanding teori.

Mata kuliah ini mencakup konsep dasar game development dan game design, pengenalan Unity Engine dan workflow pengembangan game, pemrograman C\# untuk Unity, pengelolaan game object dan komponen, sistem input dan kontrol pemain, fisika dan collision detection, animasi, UI/UX game, audio, level design, AI sederhana, optimasi, serta proses build dan deployment game. Melalui pembelajaran bertahap dan proyek akhir yang komprehensif, mahasiswa diharapkan mampu merancang, mengembangkan, menguji, dan mempublikasikan game yang fungsional dan menarik.

\section{Capaian Pembelajaran Mata Kuliah (CPMK)}

Setelah menyelesaikan mata kuliah ini, mahasiswa diharapkan mampu:

\begin{enumerate}[label=CPMK\arabic*:]
\item Menganalisis konsep dasar game development, game design, dan genre game serta menjelaskan workflow pengembangan game menggunakan Unity Engine.

\item Menerapkan pemrograman C\# untuk Unity dengan menguasai struktur skrip, fungsi dasar (\texttt{Start()}, \texttt{Update()}), variabel, kontrol alur program, dan konsep object-oriented programming dalam konteks game development.

\item Mengoperasikan Unity Editor dengan mahir dalam mengelola game object, komponen (Transform, Rigidbody, Collider), prefab, scene, dan manajemen aset (grafik, model 3D, audio).

\item Mengimplementasikan sistem kontrol pemain menggunakan Input System Unity, menggerakkan karakter, dan menangani interaksi pemain dengan game world.

\item Mengaplikasikan sistem fisika Unity (Rigidbody, Collider, Physics Material) untuk simulasi gerak, gravitasi, tumbukan, dan interaksi fisik dalam game.

\item Membuat animasi karakter dan objek menggunakan Animator Controller, mengelola transisi animasi, dan mengintegrasikan animasi dengan gameplay.

\item Merancang dan mengimplementasikan UI/UX game yang efektif meliputi HUD (Heads-Up Display), menu navigasi, sistem scoring, dan feedback visual untuk pengalaman pengguna yang optimal.

\item Mengintegrasikan audio dalam game meliputi sound effect dan background music menggunakan AudioSource, AudioListener, dan Audio Mixer untuk meningkatkan immersi gameplay.

\item Merancang level design dan game mechanics yang menarik, termasuk sistem progresi, tantangan, dan reward yang seimbang.

\item Mengimplementasikan AI sederhana untuk NPC (Non-Player Character) menggunakan teknik dasar seperti pathfinding dengan NavMesh, state machine, dan behavior scripting.

\item Menganalisis dan mengoptimalkan performa game melalui profiling, debugging, dan teknik optimasi aset serta kode untuk berbagai platform target.

\item Membangun dan mendeploy game ke berbagai platform (PC, Mobile, Web) dengan memahami proses build, pengaturan platform-specific, dan distribusi game.

\item Merancang dan mengembangkan proyek game utuh yang mencakup semua aspek pembelajaran mulai dari konsep, desain, implementasi, pengujian, hingga publikasi dengan menerapkan best practices industri game development.
\end{enumerate}

\section{Pemetaan CPMK terhadap CPL}

\begin{table}[h]
\centering
\caption{Pemetaan CPMK terhadap Capaian Pembelajaran Lulusan (CPL)}
\label{tab:pemetaan_cpl}
\small
\begin{tabular}{|p{1.5cm}|p{7cm}|p{6cm}|}
\hline
\textbf{CPMK} & \textbf{Deskripsi Singkat} & \textbf{CPL yang Terkait} \\
\hline
CPMK1 & Analisis konsep game development dan Unity workflow & \textbf{CPL1:} Sikap (Bertanggung jawab, etis, profesional) \\
 & & \textbf{CPL2:} Pengetahuan (Menguasai konsep dan prinsip) \\
\hline
CPMK2 & Penerapan pemrograman C\# untuk Unity & \textbf{CPL3:} Keterampilan Umum (Kemampuan berpikir logis) \\
 & & \textbf{CPL4:} Keterampilan Khusus (Kemampuan pemrograman) \\
\hline
CPMK3 & Operasi Unity Editor dan manajemen aset & \textbf{CPL4:} Keterampilan Khusus (Penggunaan tools) \\
\hline
CPMK4 & Implementasi sistem kontrol pemain & \textbf{CPL3:} Keterampilan Umum (Problem solving) \\
 & & \textbf{CPL4:} Keterampilan Khusus (Implementasi sistem) \\
\hline
CPMK5 & Aplikasi sistem fisika Unity & \textbf{CPL2:} Pengetahuan (Konsep fisika simulasi) \\
 & & \textbf{CPL4:} Keterampilan Khusus (Implementasi fisika) \\
\hline
CPMK6 & Pembuatan animasi karakter dan objek & \textbf{CPL4:} Keterampilan Khusus (Animasi dan visual) \\
\hline
CPMK7 & Perancangan UI/UX game & \textbf{CPL3:} Keterampilan Umum (Komunikasi visual) \\
 & & \textbf{CPL4:} Keterampilan Khusus (UI/UX design) \\
\hline
CPMK8 & Integrasi audio dalam game & \textbf{CPL4:} Keterampilan Khusus (Audio integration) \\
\hline
CPMK9 & Perancangan level design dan mechanics & \textbf{CPL2:} Pengetahuan (Game design principles) \\
 & & \textbf{CPL3:} Keterampilan Umum (Kreativitas) \\
\hline
CPMK10 & Implementasi AI sederhana & \textbf{CPL2:} Pengetahuan (Konsep AI dasar) \\
 & & \textbf{CPL4:} Keterampilan Khusus (AI implementation) \\
\hline
CPMK11 & Analisis dan optimasi performa & \textbf{CPL3:} Keterampilan Umum (Analisis kritis) \\
 & & \textbf{CPL4:} Keterampilan Khusus (Optimization) \\
\hline
CPMK12 & Build dan deployment game & \textbf{CPL3:} Keterampilan Umum (Project management) \\
 & & \textbf{CPL4:} Keterampilan Khusus (Deployment) \\
\hline
CPMK13 & Pengembangan proyek game utuh & \textbf{CPL1:} Sikap (Komitmen, profesionalisme) \\
 & & \textbf{CPL3:} Keterampilan Umum (Komprehensif) \\
 & & \textbf{CPL4:} Keterampilan Khusus (Full-stack development) \\
\hline
\end{tabular}
\end{table}

\textit{Catatan: CPL (Capaian Pembelajaran Lulusan) disesuaikan dengan program studi masing-masing.}

\section{Bahan Kajian}

\subsection{Konsep Dasar Game Development}
\begin{itemize}
\item Pengertian game, elemen-elemen game (mechanics, dynamics, aesthetics)
\item Genre game dan karakteristiknya
\item Game design document (GDD) dan konsep dasar game design
\item Lifecycle pengembangan game (pre-production, production, post-production)
\item Peran dan tanggung jawab dalam tim pengembangan game
\end{itemize}

\subsection{Pengenalan Unity Engine}
\begin{itemize}
\item Sejarah dan keunggulan Unity Engine
\item Instalasi dan setup Unity Hub
\item Unity Editor interface (Scene view, Game view, Hierarchy, Inspector, Project, Console)
\item Workflow pengembangan game di Unity
\item Unity Asset Store dan manajemen package
\end{itemize}

\subsection{Dasar Pemrograman C\# untuk Unity}
\begin{itemize}
\item Struktur skrip C\# di Unity (MonoBehaviour)
\item Fungsi \texttt{Start()}, \texttt{Update()}, \texttt{Awake()}, \texttt{OnEnable()}
\item Variabel dan tipe data (int, float, bool, string, Vector2, Vector3)
\item Kontrol alur program (if-else, switch, loop)
\item Fungsi dan method
\item Konsep object-oriented programming (class, inheritance, polymorphism)
\item Coroutine dan async programming dasar
\item Event dan delegate
\end{itemize}

\subsection{Game Object, Component, Prefab, dan Scene}
\begin{itemize}
\item Konsep GameObject dan Hierarchy
\item Transform component (position, rotation, scale)
\item Komponen-komponen Unity (Rigidbody, Collider, Renderer, AudioSource, dll)
\item Prefab dan instantiation
\item Scene management dan Scene loading
\item Tag dan Layer
\item Parenting dan child objects
\end{itemize}

\subsection{Input System dan Player Controller}
\begin{itemize}
\item Unity Input System (legacy dan new Input System)
\item Keyboard dan mouse input
\item Touch input untuk mobile
\item Gamepad/controller input
\item Player movement (2D dan 3D)
\item Camera control dan follow camera
\item Character controller vs Rigidbody movement
\end{itemize}

\subsection{Physics, Collision, dan Rigidbody}
\begin{itemize}
\item Konsep fisika dalam game
\item Rigidbody component dan propertinya
\item Collider types (Box, Sphere, Capsule, Mesh)
\item Collision detection (OnCollisionEnter, OnTriggerEnter)
\item Physics Material dan friction
\item Raycast dan Linecast
\item Joints dan constraints
\end{itemize}

\subsection{Animation dan Animator}
\begin{itemize}
\item Konsep animasi dalam game
\item Animator Controller dan state machine
\item Animation states dan transitions
\item Parameters (float, int, bool, trigger)
\item Blend trees untuk smooth transitions
\item Animation events
\item Runtime animation control melalui script
\end{itemize}

\subsection{UI/UX Game}
\begin{itemize}
\item Canvas dan Canvas Scaler
\item UI Components (Button, Text, Image, Slider, Toggle, Input Field)
\item Event System dan UI events
\item HUD (Heads-Up Display) design
\item Menu system (Main menu, Pause menu, Settings)
\item Scoring system dan display
\item Health bar dan progress indicators
\item UI Layout (Horizontal, Vertical, Grid Layout Group)
\end{itemize}

\subsection{Audio dalam Game}
\begin{itemize}
\item Konsep audio dalam game (sound effect vs background music)
\item AudioSource dan AudioListener
\item Audio Clip dan format audio
\item 2D vs 3D audio
\item Audio Mixer dan mixing
\item Dynamic audio (volume control, pitch variation)
\item Audio pooling untuk optimasi
\end{itemize}

\subsection{Level Design dan Game Mechanics}
\begin{itemize}
\item Prinsip level design (flow, pacing, challenge)
\item Spawn system dan object pooling
\item Game state management (playing, paused, game over)
\item Win/lose conditions
\item Progression system (leveling, unlocking)
\item Power-ups dan collectibles
\item Checkpoint system
\end{itemize}

\subsection{AI Sederhana untuk Game}
\begin{itemize}
\item Konsep AI dalam game
\item NPC behavior patterns (patrol, chase, attack, flee)
\item NavMesh dan pathfinding
\item State Machine untuk AI behavior
\item Line of sight detection
\item Decision making (if-else, finite state machine)
\item AI optimization techniques
\end{itemize}

\subsection{Optimization dan Debugging}
\begin{itemize}
\item Profiling tools (Unity Profiler, Frame Debugger)
\item Performance bottlenecks (CPU, GPU, memory)
\item Optimization techniques:
\begin{itemize}
\item Object pooling
\item LOD (Level of Detail)
\item Occlusion culling
\item Texture compression
\item Code optimization
\end{itemize}
\item Debugging tools dan techniques
\item Logging dan error handling
\end{itemize}

\subsection{Build dan Deployment}
\begin{itemize}
\item Build settings dan platform selection
\item Platform-specific settings (PC, Android, iOS, WebGL)
\item Asset optimization untuk target platform
\item Build process dan troubleshooting
\item Testing pada target platform
\item Distribution channels (itch.io, Google Play, App Store, Steam)
\item Version control dengan Git untuk game projects
\end{itemize}

\section{Rencana Pembelajaran Semester (RPS)}

\begin{longtable}{|p{0.8cm}|p{3.5cm}|p{4cm}|p{4cm}|p{2.5cm}|}
\hline
\textbf{No} & \textbf{Tujuan Pembelajaran} & \textbf{Materi Pembelajaran} & \textbf{Aktivitas Praktik} & \textbf{Evaluasi} \\
\hline
\endfirsthead

\hline
\textbf{No} & \textbf{Tujuan Pembelajaran} & \textbf{Materi Pembelajaran} & \textbf{Aktivitas Praktik} & \textbf{Evaluasi} \\
\hline
\endhead

\hline
\endfoot

\hline
\endlastfoot

1 & Mahasiswa mampu menjelaskan konsep dasar game development, genre game, dan workflow pengembangan game menggunakan Unity. & 
\begin{itemize}[leftmargin=*,nosep]
\item Konsep dasar game dan elemen-elemen game
\item Genre game dan karakteristiknya
\item Lifecycle pengembangan game
\item Pengenalan Unity Engine dan keunggulannya
\item Instalasi Unity Hub dan Unity Editor
\end{itemize} & 
\begin{itemize}[leftmargin=*,nosep]
\item Diskusi tentang game favorit dan analisis elemen-elemennya
\item Instalasi Unity Hub dan Unity Editor
\item Eksplorasi Unity Editor interface
\item Membuat project Unity pertama
\end{itemize} & 
Kuis: Konsep dasar game development \\
\hline

2 & Mahasiswa mampu mengoperasikan Unity Editor dengan mahir dan memahami struktur project Unity. & 
\begin{itemize}[leftmargin=*,nosep]
\item Unity Editor interface detail (Scene, Game, Hierarchy, Inspector, Project, Console)
\item Workflow pengembangan game di Unity
\item Membuat dan mengelola Scene
\item Unity Asset Store dan package management
\end{itemize} & 
\begin{itemize}[leftmargin=*,nosep]
\item Praktik navigasi Unity Editor
\item Membuat Scene baru dengan berbagai objek dasar
\item Import asset dari Asset Store
\item Setup project structure yang rapi
\end{itemize} & 
Tugas: Setup project Unity dengan struktur folder yang baik \\
\hline

3 & Mahasiswa mampu membuat skrip C\# dasar untuk Unity dan memahami lifecycle MonoBehaviour. & 
\begin{itemize}[leftmargin=*,nosep]
\item Struktur skrip C\# di Unity (MonoBehaviour)
\item Fungsi \texttt{Start()}, \texttt{Update()}, \texttt{Awake()}, \texttt{OnEnable()}
\item Variabel dan tipe data dasar
\item Debug.Log untuk debugging
\end{itemize} & 
\begin{itemize}[leftmargin=*,nosep]
\item Membuat skrip C\# pertama di Unity
\item Praktik penggunaan fungsi lifecycle
\item Manipulasi variabel dan output ke Console
\item Latihan membuat skrip sederhana (rotasi objek)
\end{itemize} & 
Tugas: Membuat skrip rotasi dan pergerakan objek sederhana \\
\hline

4 & Mahasiswa mampu menguasai konsep GameObject, Component, Transform, dan Prefab. & 
\begin{itemize}[leftmargin=*,nosep]
\item Konsep GameObject dan Hierarchy
\item Transform component (position, rotation, scale)
\item Komponen-komponen Unity dasar
\item Prefab dan instantiation
\item Tag dan Layer
\end{itemize} & 
\begin{itemize}[leftmargin=*,nosep]
\item Praktik manipulasi Transform melalui Inspector dan script
\item Membuat Prefab dari GameObject
\item Instantiate Prefab melalui script
\item Latihan membuat scene dengan berbagai objek dan prefab
\end{itemize} & 
Tugas: Membuat scene dengan prefab system \\
\hline

5 & Mahasiswa mampu mengimplementasikan sistem input dan kontrol pemain dasar. & 
\begin{itemize}[leftmargin=*,nosep]
\item Unity Input System (legacy dan new)
\item Keyboard dan mouse input
\item Player movement 2D dan 3D
\item Camera control dasar
\end{itemize} & 
\begin{itemize}[leftmargin=*,nosep]
\item Implementasi player movement dengan keyboard
\item Membuat karakter yang dapat bergerak (WASD/Arrow keys)
\item Setup camera follow untuk karakter
\item Latihan membuat kontrol pemain yang responsif
\end{itemize} & 
Tugas: Game sederhana dengan player controller (contoh: menggerakkan kotak) \\
\hline

6 & Mahasiswa mampu mengaplikasikan sistem fisika Unity untuk simulasi gerak dan tumbukan. & 
\begin{itemize}[leftmargin=*,nosep]
\item Rigidbody component dan propertinya
\item Collider types dan collision detection
\item Physics Material
\item OnCollisionEnter dan OnTriggerEnter
\end{itemize} & 
\begin{itemize}[leftmargin=*,nosep]
\item Setup Rigidbody dan Collider pada objek
\item Membuat simulasi fisika (jatuh, tumbukan)
\item Implementasi trigger untuk collectibles
\item Latihan membuat game dengan mekanik fisika (contoh: platformer sederhana)
\end{itemize} & 
Tugas: Game dengan sistem fisika (contoh: platformer atau puzzle game sederhana) \\
\hline

7 & Mahasiswa mampu membuat animasi karakter menggunakan Animator Controller. & 
\begin{itemize}[leftmargin=*,nosep]
\item Konsep animasi dalam game
\item Animator Controller dan state machine
\item Animation states dan transitions
\item Parameters (float, int, bool, trigger)
\end{itemize} & 
\begin{itemize}[leftmargin=*,nosep]
\item Membuat Animator Controller
\item Setup animation states (idle, walk, jump)
\item Membuat transitions antar states
\item Kontrol animasi melalui script
\end{itemize} & 
Tugas: Karakter dengan animasi dasar (idle, walk, jump) \\
\hline

8 & \textbf{UTS - Ujian Tengah Semester} & 
\begin{itemize}[leftmargin=*,nosep]
\item Review materi pertemuan 1-7
\item Konsep game development
\item Unity Editor dan workflow
\item Pemrograman C\# untuk Unity
\item GameObject, Component, Prefab
\item Input System dan Player Controller
\item Physics dan Collision
\item Animation dasar
\end{itemize} & 
\begin{itemize}[leftmargin=*,nosep]
\item Ujian teori (pilihan ganda dan esai)
\item Presentasi progress proyek game individu
\item Demo game sederhana yang telah dibuat
\end{itemize} & 
UTS: Ujian teori dan presentasi proyek \\
\hline

9 & Mahasiswa mampu merancang dan mengimplementasikan UI/UX game yang efektif. & 
\begin{itemize}[leftmargin=*,nosep]
\item Canvas dan Canvas Scaler
\item UI Components (Button, Text, Image, Slider)
\item Event System dan UI events
\item HUD design (health bar, score display)
\item Menu system (Main menu, Pause menu)
\end{itemize} & 
\begin{itemize}[leftmargin=*,nosep]
\item Membuat Canvas dan setup UI
\item Implementasi HUD (score, health bar)
\item Membuat Main Menu dengan navigasi
\item Implementasi Pause Menu
\item Latihan membuat UI yang responsif
\end{itemize} & 
Tugas: Game dengan UI lengkap (HUD + Menu) \\
\hline

10 & Mahasiswa mampu mengintegrasikan audio dalam game untuk meningkatkan immersi. & 
\begin{itemize}[leftmargin=*,nosep]
\item AudioSource dan AudioListener
\item Sound effect vs background music
\item 2D vs 3D audio
\item Audio Mixer dasar
\item Dynamic audio control
\end{itemize} & 
\begin{itemize}[leftmargin=*,nosep]
\item Import audio assets
\item Setup AudioSource untuk sound effect
\item Implementasi background music
\item Kontrol volume melalui UI
\item Latihan membuat audio yang sesuai dengan gameplay
\end{itemize} & 
Tugas: Menambahkan audio (SFX + BGM) ke game \\
\hline

11 & Mahasiswa mampu merancang level design dan game mechanics yang menarik. & 
\begin{itemize}[leftmargin=*,nosep]
\item Prinsip level design (flow, pacing, challenge)
\item Game state management
\item Win/lose conditions
\item Spawn system dan object pooling
\item Collectibles dan power-ups
\end{itemize} & 
\begin{itemize}[leftmargin=*,nosep]
\item Merancang level dengan prinsip level design
\item Implementasi game state (playing, paused, game over)
\item Membuat sistem spawn untuk enemies/objects
\item Implementasi collectibles dan scoring
\item Latihan membuat level yang seimbang dan menarik
\end{itemize} & 
Tugas: Level design dengan mechanics lengkap \\
\hline

12 & Mahasiswa mampu mengimplementasikan AI sederhana untuk NPC menggunakan teknik dasar. & 
\begin{itemize}[leftmargin=*,nosep]
\item Konsep AI dalam game
\item NPC behavior patterns (patrol, chase, attack)
\item NavMesh dan pathfinding
\item State Machine untuk AI
\item Line of sight detection
\end{itemize} & 
\begin{itemize}[leftmargin=*,nosep]
\item Setup NavMesh untuk pathfinding
\item Membuat NPC dengan behavior patrol
\item Implementasi chase behavior untuk enemy
\item State Machine untuk AI (idle, patrol, chase, attack)
\item Latihan membuat AI yang menantang namun fair
\end{itemize} & 
Tugas: Game dengan AI enemy (patrol + chase) \\
\hline

13 & Mahasiswa mampu menganalisis dan mengoptimalkan performa game. & 
\begin{itemize}[leftmargin=*,nosep]
\item Unity Profiler dan Frame Debugger
\item Performance bottlenecks (CPU, GPU, memory)
\item Optimization techniques (object pooling, LOD, occlusion culling)
\item Debugging tools dan techniques
\end{itemize} & 
\begin{itemize}[leftmargin=*,nosep]
\item Menggunakan Profiler untuk analisis performa
\item Implementasi object pooling
\item Optimasi texture dan model
\item Debugging game yang sudah dibuat
\item Latihan mengoptimalkan game untuk performa lebih baik
\end{itemize} & 
Tugas: Optimasi game yang telah dibuat \\
\hline

14 & Mahasiswa mampu membangun dan mendeploy game ke berbagai platform. & 
\begin{itemize}[leftmargin=*,nosep]
\item Build settings dan platform selection
\item Platform-specific settings (PC, Android, WebGL)
\item Asset optimization untuk target platform
\item Build process dan troubleshooting
\item Distribution channels
\end{itemize} & 
\begin{itemize}[leftmargin=*,nosep]
\item Setup build settings untuk PC
\item Build game untuk WebGL
\item Setup build untuk Android (jika memungkinkan)
\item Testing build pada target platform
\item Upload ke platform distribusi (itch.io)
\end{itemize} & 
Tugas: Build dan deploy game ke platform \\
\hline

15 & Mahasiswa mampu mengembangkan proyek game utuh dengan menerapkan semua aspek pembelajaran. & 
\begin{itemize}[leftmargin=*,nosep]
\item Review semua materi pembelajaran
\item Best practices game development
\item Project management untuk game development
\item Version control dengan Git
\item Testing dan quality assurance
\end{itemize} & 
\begin{itemize}[leftmargin=*,nosep]
\item Finalisasi proyek game akhir
\item Implementasi fitur-fitur yang belum ada
\item Testing menyeluruh
\item Dokumentasi proyek
\item Persiapan presentasi final
\end{itemize} & 
Progress check: Review proyek akhir \\
\hline

16 & \textbf{UAS - Ujian Akhir Semester} & 
\begin{itemize}[leftmargin=*,nosep]
\item Presentasi proyek game akhir
\item Demo game yang telah dikembangkan
\item Q\&A tentang implementasi dan konsep
\item Refleksi pembelajaran
\end{itemize} & 
\begin{itemize}[leftmargin=*,nosep]
\item Presentasi proyek game akhir (15-20 menit)
\item Demo gameplay
\item Penjelasan implementasi teknis
\item Diskusi dan evaluasi
\end{itemize} & 
UAS: Presentasi dan demo proyek akhir \\
\hline

\end{longtable}

\section{Metode Pembelajaran}

Mata kuliah ini menggunakan pendekatan pembelajaran yang berorientasi pada praktik dengan metode sebagai berikut:

\subsection{Praktikum}
Setiap pertemuan mencakup sesi praktikum langsung menggunakan Unity Editor dan Visual Studio/Visual Studio Code. Mahasiswa akan melakukan hands-on practice dengan bimbingan dosen untuk menguasai tools dan teknik pengembangan game.

\subsection{Project-Based Learning (PBL)}
Mahasiswa akan mengembangkan proyek game secara bertahap mulai dari pertemuan pertama hingga akhir semester. Proyek ini akan berkembang seiring dengan materi yang dipelajari, sehingga mahasiswa dapat langsung menerapkan konsep yang dipelajari ke dalam proyek nyata.

\subsection{Problem-Based Learning}
Mahasiswa diberikan studi kasus dan permasalahan nyata dalam pengembangan game, kemudian diminta untuk menganalisis dan mencari solusi menggunakan tools dan teknik yang telah dipelajari.

\subsection{Studi Kasus}
Analisis game-game populer dan mekaniknya untuk memahami implementasi konsep-konsep game development dalam produk komersial.

\subsection{Presentasi dan Diskusi}
Mahasiswa mempresentasikan progress proyek mereka secara berkala dan berdiskusi tentang tantangan serta solusi yang ditemukan dalam proses pengembangan.

\subsection{Peer Review}
Mahasiswa saling memberikan feedback terhadap proyek teman sekelas untuk meningkatkan kualitas dan pembelajaran kolaboratif.

\section{Evaluasi dan Penilaian}

\subsection{Bobot Penilaian}

\begin{table}[h]
\centering
\caption{Distribusi Bobot Penilaian}
\label{tab:bobot_penilaian}
\begin{tabular}{|l|c|}
\hline
\textbf{Komponen Penilaian} & \textbf{Bobot (\%)} \\
\hline
Tugas Praktikum (8 tugas) & 25\% \\
Kuis (2 kuis) & 10\% \\
UTS (Ujian Tengah Semester) & 20\% \\
Proyek Akhir & 30\% \\
UAS (Ujian Akhir Semester) & 15\% \\
\hline
\textbf{Total} & \textbf{100\%} \\
\hline
\end{tabular}
\end{table}

\subsection{Rubrik Penilaian Berbasis OBE}

\subsubsection{Rubrik Penilaian Tugas Praktikum}

\begin{table}[h]
\centering
\caption{Rubrik Penilaian Tugas Praktikum}
\label{tab:rubrik_tugas}
\small
\begin{tabular}{|p{2cm}|p{3.5cm}|p{3.5cm}|p{3.5cm}|}
\hline
\textbf{Kriteria} & \textbf{Sangat Baik (85-100)} & \textbf{Baik (70-84)} & \textbf{Cukup (60-69)} \\
\hline
\textbf{Fungsionalitas} & Semua fitur berfungsi dengan sempurna, tidak ada bug yang mengganggu gameplay & Sebagian besar fitur berfungsi, ada beberapa bug minor & Fitur dasar berfungsi, namun ada bug yang mempengaruhi gameplay \\
\hline
\textbf{Implementasi Teknis} & Kode rapi, terstruktur, menggunakan best practices, komentar jelas & Kode cukup rapi, struktur baik, beberapa best practices diterapkan & Kode berfungsi namun kurang rapi, struktur perlu perbaikan \\
\hline
\textbf{Kreativitas} & Implementasi kreatif dan inovatif, menambahkan fitur tambahan yang relevan & Implementasi standar dengan beberapa variasi kreatif & Implementasi sesuai instruksi tanpa variasi \\
\hline
\textbf{Ketepatan Waktu} & Dikumpulkan tepat waktu atau lebih awal & Dikumpulkan dengan keterlambatan maksimal 1 hari & Dikumpulkan dengan keterlambatan lebih dari 1 hari \\
\hline
\end{tabular}
\end{table}

\subsubsection{Rubrik Penilaian Proyek Akhir}

\begin{table}[h]
\centering
\caption{Rubrik Penilaian Proyek Akhir}
\label{tab:rubrik_proyek}
\small
\begin{tabular}{|p{2cm}|p{3.5cm}|p{3.5cm}|p{3.5cm}|}
\hline
\textbf{Kriteria} & \textbf{Sangat Baik (85-100)} & \textbf{Baik (70-84)} & \textbf{Cukup (60-69)} \\
\hline
\textbf{Kompleksitas dan Scope} & Game lengkap dengan multiple levels, berbagai mechanics, dan fitur advanced & Game cukup lengkap dengan beberapa levels dan mechanics yang bervariasi & Game dasar dengan 1-2 levels dan mechanics sederhana \\
\hline
\textbf{Fungsionalitas} & Semua sistem berfungsi sempurna, gameplay smooth, tidak ada bug kritis & Sebagian besar sistem berfungsi, ada beberapa bug minor & Fitur dasar berfungsi, namun ada bug yang mempengaruhi experience \\
\hline
\textbf{Desain dan UX} & UI/UX sangat menarik dan intuitif, visual polish tinggi, audio terintegrasi dengan baik & UI/UX baik dan fungsional, visual cukup menarik, audio ada namun kurang optimal & UI/UX dasar, visual sederhana, audio minimal atau tidak ada \\
\hline
\textbf{Kualitas Kode} & Kode sangat rapi, terstruktur dengan baik, menggunakan design patterns, dokumentasi lengkap & Kode rapi dan terstruktur, beberapa design patterns digunakan, dokumentasi ada & Kode berfungsi namun perlu refactoring, dokumentasi minimal \\
\hline
\textbf{Presentasi} & Presentasi sangat jelas, menjelaskan semua aspek dengan detail, demo smooth & Presentasi jelas, menjelaskan aspek utama, demo berjalan dengan baik & Presentasi cukup, beberapa aspek kurang jelas, demo ada beberapa masalah \\
\hline
\textbf{Inovasi dan Kreativitas} & Game sangat kreatif dengan konsep unik, inovasi dalam mechanics atau implementasi & Game kreatif dengan beberapa elemen unik & Game mengikuti konsep umum tanpa banyak inovasi \\
\hline
\end{tabular}
\end{table}

\subsubsection{Rubrik Penilaian Ujian (UTS/UAS)}

\begin{table}[h]
\centering
\caption{Rubrik Penilaian Ujian Teori}
\label{tab:rubrik_ujian}
\small
\begin{tabular}{|p{2cm}|p{3.5cm}|p{3.5cm}|p{3.5cm}|}
\hline
\textbf{Kriteria} & \textbf{Sangat Baik (85-100)} & \textbf{Baik (70-84)} & \textbf{Cukup (60-69)} \\
\hline
\textbf{Pemahaman Konsep} & Menjawab dengan akurat, menunjukkan pemahaman mendalam, memberikan contoh yang relevan & Menjawab dengan benar, menunjukkan pemahaman baik, memberikan beberapa contoh & Menjawab sebagian benar, pemahaman dasar ada namun kurang mendalam \\
\hline
\textbf{Aplikasi Praktis} & Mampu mengaplikasikan konsep ke situasi nyata dengan solusi yang tepat dan kreatif & Mampu mengaplikasikan konsep dengan solusi yang tepat & Mampu mengaplikasikan konsep dasar namun solusi kurang optimal \\
\hline
\textbf{Analisis dan Evaluasi} & Analisis mendalam, evaluasi kritis, memberikan insight yang berharga & Analisis cukup baik, evaluasi ada namun kurang mendalam & Analisis dasar, evaluasi minimal \\
\hline
\end{tabular}
\end{table}

\subsection{Strategi Penilaian OBE}

Penilaian dalam mata kuliah ini mengikuti prinsip Outcome-Based Education dengan fokus pada:

\begin{enumerate}
\item \textbf{Authentic Assessment}: Penilaian dilakukan melalui proyek nyata pengembangan game yang mencerminkan kompetensi yang diharapkan.

\item \textbf{Continuous Assessment}: Penilaian dilakukan secara berkelanjutan melalui tugas praktikum, bukan hanya pada ujian akhir.

\item \textbf{Multiple Assessment Methods}: Menggunakan berbagai metode penilaian (tugas praktik, proyek, ujian teori, presentasi) untuk menilai berbagai aspek kompetensi.

\item \textbf{Competency-Based}: Penilaian berdasarkan pencapaian kompetensi (CPMK) yang telah ditetapkan, bukan hanya berdasarkan nilai akhir.

\item \textbf{Feedback-Oriented}: Setiap penilaian disertai dengan feedback konstruktif untuk membantu mahasiswa meningkatkan kompetensinya.
\end{enumerate}

\section{Media dan Perangkat Pembelajaran}

\subsection{Perangkat Lunak}
\begin{itemize}
\item \textbf{Unity Hub}: Untuk manajemen instalasi Unity dan project
\item \textbf{Unity Editor} (versi LTS terbaru): Game engine utama
\item \textbf{Visual Studio} atau \textbf{Visual Studio Code}: IDE untuk pemrograman C\#
\item \textbf{Git}: Version control system
\item \textbf{Unity Asset Store}: Sumber asset gratis dan berbayar
\item \textbf{Audacity} atau \textbf{Adobe Audition}: Editor audio (opsional)
\item \textbf{GIMP} atau \textbf{Photoshop}: Editor gambar (opsional)
\end{itemize}

\subsection{Perangkat Keras}
\begin{itemize}
\item Komputer/laptop dengan spesifikasi minimal:
\begin{itemize}
\item Processor: Intel Core i5 atau AMD equivalent
\item RAM: 8 GB (16 GB direkomendasikan)
\item Graphics Card: DirectX 11 compatible
\item Storage: 20 GB free space untuk Unity dan project
\end{itemize}
\item Mouse dengan scroll wheel (penting untuk navigasi Unity Editor)
\item Headphone/speaker untuk audio development
\end{itemize}

\subsection{Sumber Belajar Online}
\begin{itemize}
\item \textbf{Unity Learn}: Platform pembelajaran resmi Unity (https://learn.unity.com/)
\item \textbf{Unity Documentation}: Dokumentasi resmi Unity (https://docs.unity3d.com/)
\item \textbf{Brackeys YouTube Channel}: Tutorial Unity populer
\item \textbf{Code Monkey YouTube Channel}: Tutorial Unity advanced
\item \textbf{Unity Forum}: Komunitas dan diskusi (https://forum.unity.com/)
\item \textbf{Stack Overflow}: Q\&A untuk troubleshooting
\end{itemize}

\subsection{Asset dan Resources}
\begin{itemize}
\item Unity Asset Store (free assets)
\item Kenney.nl (free game assets)
\item OpenGameArt.org (free game assets)
\item Freesound.org (free audio assets)
\item Unity Learn Project Files
\end{itemize}

\section{Referensi}

\subsection{Referensi Utama}

\begin{enumerate}
\item Unity Technologies. (2024). \textit{Unity User Manual}. Unity Documentation. https://docs.unity3d.com/Manual/index.html

\item Unity Technologies. (2024). \textit{Unity Scripting API}. Unity Documentation. https://docs.unity3d.com/ScriptReference/index.html

\item Microsoft. (2024). \textit{C\# Documentation}. Microsoft Learn. https://learn.microsoft.com/en-us/dotnet/csharp/

\item Schell, J. (2019). \textit{The Art of Game Design: A Book of Lenses} (3rd ed.). CRC Press.

\item Nystrom, R. (2014). \textit{Game Programming Patterns}. Genever Benning. https://gameprogrammingpatterns.com/

\item Gregory, J. (2018). \textit{Game Engine Architecture} (3rd ed.). CRC Press.

\item Salen, K., \& Zimmerman, E. (2004). \textit{Rules of Play: Game Design Fundamentals}. MIT Press.
\end{enumerate}

\subsection{Referensi Pendukung}

\begin{enumerate}
\item Unity Learn. (2024). \textit{Unity Learn Platform}. https://learn.unity.com/

\item Brackeys. (2024). \textit{Unity Tutorials}. YouTube Channel. https://www.youtube.com/c/Brackeys

\item Code Monkey. (2024). \textit{Unity Tutorials}. YouTube Channel. https://www.youtube.com/c/CodeMonkeyUnity

\item GameDev.tv. (2024). \textit{Complete C\# Unity Game Developer 3D}. Udemy Course.

\item Fogue, M. (2023). \textit{Unity Game Development Cookbook}. O'Reilly Media.

\item Hocking, J. (2018). \textit{Unity in Action: Multiplatform Game Development in C\#} (2nd ed.). Manning Publications.

\item Thorn, A. (2017). \textit{Learn to Code by Making Games: The Complete Unity Developer}. Udemy Course.

\item Gibson, J. (2014). \textit{Introduction to Game Design, Prototyping, and Development: From Concept to Playable Game with Unity and C\#}. Addison-Wesley Professional.

\item Blackman, S. (2013). \textit{Beginning 3D Game Development with Unity 4: All-in-One, Multi-Platform Game Development}. Apress.

\item Smith, M. (2020). \textit{Unity Game Development in 24 Hours}. Sams Publishing.
\end{enumerate}

\subsection{Referensi Online dan Komunitas}

\begin{itemize}
\item Unity Forum: https://forum.unity.com/
\item Unity Answers: https://answers.unity.com/
\item Stack Overflow (Unity tag): https://stackoverflow.com/questions/tagged/unity3d
\item Reddit r/Unity3D: https://www.reddit.com/r/Unity3D/
\item Unity Discord Server
\item GameDev.net: https://www.gamedev.net/
\end{itemize}

\vspace{2cm}

\begin{center}
\textit{Dokumen ini disusun berdasarkan kurikulum Outcome-Based Education (OBE) dan akan direview secara berkala untuk memastikan relevansi dengan kebutuhan industri dan perkembangan teknologi game development.}
\end{center}

\end{document}
